\section{Related Works}

\subsection{World Models Design}
World models characterize how environment states evolve over time under given interaction conditions, thereby providing effective support for tasks such as counterfactual simulation, planning, and decision-making \cite{ha2018world}.
Early world models primarily relied on multimodal large language models (MLLMs) \cite{hurst2024gpt,Qwen3-VL,bai2025univg,chu2025gpg} that represent world states through textual abstractions \cite{yang2024rila,shridhar2020alfworld}.  Recent advances in video generation \cite{sora_initial,wan2025,mao2025omni,wu2026latent,zhu2026artifact} have driven a shift toward video-based world models, which offer a more expressive and grounded representation of complex environments and have emerged as a dominant paradigm in the field \cite{ding2025understanding,zhu2024sora,zheng2026code2world}. In this work, we focus on world models built upon video generation.


Across different application domains, video-based world models have followed distinct yet intrinsically related technical trajectories. In autonomous driving, world models primarily focus on long-horizon traffic scene evolution and the decision-making of vehicle agents \cite{feng2025survey}. Representative works such as GAIA \cite{hu2023gaia}, DriveDreamer \cite{wang2024drivedreamer}, DrivingWorld \cite{hu2024drivingworld}, and Vista \cite{gao2024vista} leverage action-conditioned future frame prediction to support planning and simulation. 
In embodied intelligence and robotics, world models place greater emphasis on object-centric dynamics and manipulation control \cite{long2025survey}. Methods such as IRASim \cite{zhu2024irasim}, Cosmos \cite{agarwal2025Cosmos}, RoboScape \cite{shang2025roboscape} and LVP \cite{chen2025large} tightly integrate perception, action, and physical reasoning to simulate interaction-driven environment changes. 
In game environments, works including Genie \cite{bruce2024genie,parker2024genie}, Matrix-Game \cite{zhang2025matrix,he2025matrix}, WorldPlay \cite{sun2025worldplay}, and Hunyuan-GameCraft \cite{li2025hunyuan,tang2025hunyuan} aim to construct highly interactive and playable virtual worlds. Despite differences in input modalities, action spaces, and domain-specific constraints, these methods share a common objective: learning how the environment responds coherently to different interaction instructions. This highlights interaction as a core capability of world modeling \cite{ding2025understanding,agarwal2025Cosmos}.
Motivated by this, our benchmark takes interaction as the central axis for evaluating world models.

\subsection{World Models Evaluation}

Despite the rapid progress of video-based world models, the development of corresponding evaluation benchmarks has remained relatively limited \cite{ding2025understanding}. Early studies \cite{feng2024matrix,yu2025gamefactory,guo2025mineworld,chen2025finger,li2025next} primarily rely on generic metrics, such as FID \cite{heusel2017gans}, IS \cite{salimans2016improved}, FVD \cite{unterthiner2019fvd}, which often exhibit significant deviations from human perceptual judgments \cite{ding2022cogview2,otani2023toward, ling2025vmbench, wu2026imagerysearch}.
Subsequently, several evaluation tools originally designed for video generation, such as VBench \cite{huang2024vbench}, have been introduced \cite{che2024gamegen,tang2025hunyuan,ling2025vmbench,feng2025narrlv,zhu2026artifact}. While these benchmarks play an important role in assessing overall visual quality and text–video alignment \cite{liu2024survey,huang2024mmgenbench}, they struggle to adequately characterize the core interactive capabilities of world model tasks. As a result, such metrics provide only limited insights for the design and analysis of interactive world models.
Moreover, WorldScore \cite{duan2025worldscore} has been proposed as a benchmark specifically tailored to world models. It focuses on evaluating a model’s ability to generate geometrically consistent 3D scenes under viewpoint changes, emphasizing spatial coherence and geometric realism. Although this represents an important step toward world-model-aware evaluation, the considered form of interaction is largely restricted to camera motion. In contrast, contemporary world models increasingly emphasize a broader range of interaction types \cite{tang2025hunyuan,chen2025large}.
Motivated by this gap, we introduce Omni-WorldBench, an interaction-centric evaluation benchmark that systematically covers multiple levels of interaction complexity. We hope that Omni-WorldBench can serve as a comprehensive tool for characterizing the interactive expressiveness of world models.